% ==============================================================================
% ПАРТИЯ B03-005: Глава 1: Введение и рисунок 1 (схемы БАВР, Visio OLE)
% Диапазон: DISS-005-B0207 .. DISS-005-B0235 (30 блоков, включая B0206 в заголовке главы)
% ==============================================================================

% DISS-005-B0207
\section{Введение}\label{sec:ch1_intro}

% DISS-005-B0208
Необходимость обеспечения надежного электроснабжения для непрерывности технологических процессов обуславливает активное применение устройств быстродействующего автоматического ввода резерва (БАВР) на промышленных предприятиях. БАВР являются усовершенствованием традиционных устройств автоматического ввода резерва (АВР). В большинстве устройств АВР используются пусковые органы напряжения. Их уставки по времени срабатывания требуется согласовывать с временами срабатывания защит отходящих присоединений и другими автоматическими устройствами. Это может увеличить время переключения (и время нарушения электроснабжения) до нескольких секунд [1,2].\par\vspace{0.2cm}

% DISS-005-B0209
Устройства БАВР используют селективные пусковые, разрешающие и блокирующие органы. Это позволяет определять место (зону) аварии без необходимости отстройки по времени от других защит элементов сети, обеспечивая высокую чувствительность.\par\vspace{0.2cm}

% DISS-005-B0210
Основные пусковые органы БАВР [3-9] – контроль минимального напряжения (ANSI код 27/27S) и пуск по углу между векторами напряжения (для определения торможения двигательной нагрузки). Устройства некоторых производителей, преимущественно иностранных, в качестве основного пускового органа используют не «пуск по углу», а сигнал частоты или его изменение (ANSI код 81L) [8, 9], однако его физическая суть от этого не меняется.\par\vspace{0.2cm}

% DISS-005-B0211
Основной разрешающий орган БАВР – контроль повышенного напряжения (ANSI 59), который контролирует достаточность напряжения на резервном источнике.\par\vspace{0.2cm}

% DISS-005-B0212
Основные блокирующие органы БАВР – это токовая защита от междуфазных коротких замыканий (КЗ) (ANSI код 50/51), реле направления мощности (РНМ) и контроль напряжения нулевой последовательности (ANSI код 59N).\par\vspace{0.2cm}

% DISS-005-B0213
Безусловно, используются и другие органы для обеспечения функционирования в различных режимах, например:\par\vspace{0.2cm}

\begin{itemize}
  % DISS-005-B0214
  \item Защита минимального тока (ANSI код 37) используется для исключения некорректного срабатывания при переводе секции в холостой ход;
  % DISS-005-B0215
  \item Блокировка по сумме токов применяется для исключения срабатывания в режимах, когда ток превышает максимально допустимый для одного из вводов, что может возникать на несимметричных подстанциях;
  % DISS-005-B0216
  \item Защиты по току нулевой последовательности (ANSI код 50N/51N и 67N) могут выступать в качестве пусковых. Это актуально для объектов, где ОЗЗ является частым случаем и перерастает в двухфазное замыкание на землю. Данный вопрос будет рассмотрен отдельно в главе ХХХХ\commentnote{DISS-005-C0007}{Евдаков Алексей}{Вставить ссылку}.
  % DISS-005-B0217
  \item Контроль синхронизма иногда реализуется отдельным органом (ANSI код 25). Использование данной функции актуально при подключении мощной синхронной двигательной нагрузки. Однако при исправном и корректном функционировании, правильность работы может быть обеспечена базовым набором органов.
\end{itemize}

% DISS-005-B0218
Но эти функции не являются минимально требуемыми для корректного функционирования, и встречаются лишь на ряде особых объектов. Базовая же минимальная схема подключения и схематическое изображение упрощенной логики БАВР приведено на рисунке 1 (подписи в сигналах «B1/B2» указывают на номер секции и соответствующий ему выключатель).\par\vspace{0.3cm}

% DISS-005-B0219 .. DISS-005-B0224 (Рисунок 1)
\begin{figure}[htbp]
\centering
\begin{minipage}{0.48\textwidth}
  \centering
  \includegraphics[width=\linewidth]{image1.png}
  \par\vspace{1mm}
  \textbf{а)}
  \commentnote{DISS-005-C0008}{Евдаков Алексей}{Перевести на русский}
\end{minipage}\hfill
\begin{minipage}{0.48\textwidth}
  \centering
  \includegraphics[width=\linewidth]{image2.png}
  \par\vspace{1mm}
  \textbf{б)}
\end{minipage}
\caption[Схемы а) подключения БАВР; б) логики срабатывания БАВР]{Схемы а) подключения БАВР; б) логики срабатывания БАВР\wordchoice{}{.}\wordrevnote{DISS-005-REV0001}{Евдаков Алексей}{2025-01-27}{вставка точки}}
\label{fig:bavr_schemes}
\end{figure}

% DISS-005-B0225
Использование этих базовых органов в алгоритмах БАВР обеспечивает практически мгновенную реакцию на аварию, что позволяет обеспечить [10, 11]:\par\vspace{0.2cm}

\begin{itemize}
  % DISS-005-B0226
  \item обеспечить надежное и непрерывное электроснабжение потребителей в случае аварийных и ненормальных режимов в питающих и распределительных электрических сетях;
  % DISS-005-B0227
  \item обеспечить успешный самозапуск электродвигателей;
  % DISS-005-B0228
  \item предотвратить дорогостоящие перерывы в работе технологического оборудования;
  % DISS-005-B0229
  \item повысить ресурс электродвигателей, насосов, трансформаторов и другого оборудования ввиду снижения токов самозапуска в 2–3 раза;
  % DISS-005-B0230
  \item повысить уровень автоматизации производства и увеличить производительность;
  % DISS-005-B0231
  \item окупить расходы на установку комплекса БАВР приблизительно за 1 год (среднестатистическое значения).
\end{itemize}

% DISS-005-B0232
За счет повышения надежности электроснабжения и отсутствия простоя технологического оборудования, установка устройств БАВР имеет значительный экономический эффект для промышленных предприятий\par\vspace{0.2cm}

% DISS-005-B0233
Время реакции алгоритмов БАВР на аварию, которые реализованы различными производителями, составляет порядка 5÷12 мс [11]., однако указанное значение является лишь паспортным при идеальных условиях, и в реальных случаях могут наблюдаться намного большие значения [12], например, при мощной двигательной нагрузке. По данным моделирования и эксплуатации в действующих электроустановках современные устройства второго поколения производства ООО «АПС» за счёт адаптивных органов работают стабильнее, а их время реакции составляет 4÷10 мс [4]. Быстродействие АВР критично для различных ответственных нагрузок, например, для двигателей с частотным приводом, станков, осуществляющих намотку проволоки или тканей и т.д. Ввиду таких строгих требований по времени реакции на аварийный процесс, устройства БАВР всегда работают в переходных режимах. В следствии этого, важно изучение поведение различных измерительных органов на начальных временных участках аварий и коммутаций, происходящих в сети.\par\vspace{0.2cm}

% DISS-005-B0234
\noindent\textit{Добавить небольшой логический переход от этой актуальности исследования к Статистики}\par\vspace{0.3cm}

% DISS-005-B0235 (пустая строка/разделитель)
\par\vspace{0.3cm}